\section{Discussion}

\subsection{Chemical Slip as a diagnostic for LLM chemical reasoning}

Chemical Slip---the mismatch between an LLM-planned adsorption site and the site realized after physical relaxation---was introduced as a feedback signal for the closed-loop agent.
However, the slip statistics generated across 135 benchmark runs reveal a secondary value: slip frequency, stratified by surface type and adsorbate, serves as a physically grounded proxy for LLM chemical reasoning quality that is independent of any downstream task metric.

Across all three backends, the one-shot slip rate on intermetallic surfaces (71--73\%) far exceeds the rate on monometallic surfaces (20--40\%), confirming that LLM chemical priors are substantially less reliable on complex alloy compositions.
This pattern is consistent across all four backends (Gemini, Grok-4, GPT-5.4, and Claude), suggesting that the bias is shared across model families rather than being an artifact of any single training pipeline.
In contrast to text-based chemistry benchmarks (e.g., ChemBench) that test factual recall, Chemical Slip evaluates \textit{applied} reasoning---whether the model's structural prediction survives physical simulation---and may therefore complement existing LLM evaluation frameworks for scientific domains.

\subsection{Closed-loop versus open-loop: complementary paradigms}

The ablation results do not support a blanket claim that closed-loop search is superior to single-shot prediction.
On four of the five ablation cases, Gemini's full agent and single-shot baseline produce identical or near-identical energies; the closed loop adds value only on the hardest case.
This is not a weakness of the architecture but a reflection of where the need lies: when the base planner already identifies the correct site family in one shot, iterative refinement has nothing to improve.

The practical implication is that closed-loop and open-loop methods occupy different niches.
For high-throughput catalyst screening---thousands of candidates for initial down-selection---open-loop agents such as Adsorb-Agent are more efficient because they are fully parallelizable and require only a single evaluation per system.
For subsequent deep exploration of promising but poorly characterized candidates (intermetallics, multimetallics, novel compositions), closed-loop refinement adds measurable value by correcting the planner's initial biases through physical feedback.
A two-stage pipeline---open-loop screening followed by closed-loop refinement on top candidates---could combine the throughput of the former with the depth of the latter; we leave this integration to future work.

\subsection{MLFF as surrogate: scope and failure modes}

All experiments in this study use MACE-MP-0 (small) as the physical backend.
The agent's value is orthogonal to MLFF accuracy: it optimizes search strategy, not the energy function itself.
Inter-configuration energy differences are primarily determined by local adsorbate--surface bonding geometry, which equivariant graph neural networks describe faithfully.
Systematic bias in absolute energies approximately cancels across configurations on the same surface and does not affect rankings.
DFT validation (Section~3.4) tests this assumption directly.

Known MLFF failure modes---dispersion-dominated binding, strongly correlated $f$-electron systems, strong charge-transfer interfaces---define the current scope boundary of the framework.
Extending to these regimes will require either higher-fidelity MLFFs or hybrid MLFF--DFT workflows, but does not require changes to the agent architecture.

\subsection{Limitations and negative results}

\paragraph{Statistical power.}
The 5-case ablation set limits the power of nonparametric tests: no pairwise comparison reaches significance after multiple-testing correction.
We therefore interpret the ablation primarily through effect sizes and bootstrap confidence intervals, treating the consistent direction of effects across three independent LLM backends as a form of conceptual replication rather than relying on $p$-values from any single backend.
Expanding the ablation to a larger and more diverse case set is a natural next step.

\paragraph{FORBID over-pruning.}
On Gemini, GPT-5.4, and Claude, removing the FORBID constraint \textit{improves} the final energy on case~19 (by 0.201, 0.136, and 0.136~eV, respectively).
This negative result, replicated independently across three of the four backends, shows that explicit negative constraints can over-prune the search space on reactive surfaces where the globally optimal site was initially rejected.
Adaptive or soft-constraint variants of the FORBID mechanism merit further investigation.

\paragraph{Scope.}
The current implementation handles flat and stepped crystalline surfaces with a single adsorbate species.
Extensions to defective surfaces, nanoparticles, co-adsorption, and solvated interfaces are architecturally straightforward but require additional validation.

\paragraph{Platform sensitivity.}
A cross-machine control experiment (re-running Gemini and Grok-4 one-shot cases on a separate EPFL workstation) reveals that case~14 shifts by approximately 0.19--0.21~eV between the two hardware environments, while case~09 is stable within 0.003~eV.
Because both backends exhibit the same shift while preserving identical search trajectories, this is best interpreted as a platform/MLFF-environment effect rather than an LLM decision-path divergence.
Within-backend ablation comparisons (which use a single machine throughout) are unaffected, but absolute cross-backend energy rankings should be interpreted with this caveat in mind.

\paragraph{LLM reproducibility.}
LLM APIs are black-box services subject to unannounced model updates.
Even at temperature~0.0, outputs may vary across API versions.
We mitigate this by recording exact model version identifiers, providing complete prompt templates (SI-1), and releasing full trajectory logs (SI-2) so that all reported results can be independently verified from the saved data.
The four-backend comparison further demonstrates that conclusions hold across different LLM providers, reducing dependence on any single model.

\section{Conclusion}

We have presented AdsMind, a closed-loop multi-agent framework that integrates LLM-driven planning with iterative MLFF relaxation feedback for autonomous adsorption configuration search.
Three mechanisms---Chemical Slip detection, FORBID constraint memory, and autonomous termination---enable the agent to detect its own reasoning errors, avoid repeating them, and stop when further exploration is unproductive.

On a 20-surface benchmark, AdsMind achieves a 100\% valid-structure success rate on Gemini and Grok-4 (18/20 on Claude and GPT-5.4 raw, with NNH dissociation as the sole failure mode; 19/20 retry-corrected on GPT-5.4).
The controlled 5-case ablation across four LLM backends (100 runs total, all successful) reveals that:
(i) the closed loop improves search quality primarily on difficult intermetallic cases where LLM priors are unreliable;
(ii) slip feedback and FORBID constraints matter selectively---they are most valuable when the base planner struggles, but FORBID can occasionally over-constrain the search;
(iii) termination is primarily a cost-control mechanism that saves substantial compute without sacrificing energy quality; and
(iv) the iterative architecture drives heterogeneous LLM backends toward the same physical solution, reducing the four-backend energy range from 0.426~eV (one-shot) to 0.053~eV (full search), an 8.0$\times$ reduction in backend spread.

Chemical Slip, introduced here as a feedback signal, additionally serves as a physically grounded diagnostic for LLM chemical reasoning quality---a standalone methodological contribution that may extend beyond the specific application in this work.

Code, frozen experiment configurations, structured result files, and complete agent logs are available at \url{https://github.com/AI4QC/AdsMind}.
